\documentclass[10pt,a4paper]{amsart}
\usepackage[margin=19mm]{geometry}
\usepackage{amsmath,amssymb,mathtools}
\usepackage[scr=rsfs,cal=euler]{mathalfa}
\usepackage{xfrac}
\usepackage[unicode,hidelinks,bookmarks=false]{hyperref}
\newcommand{\ie}{i.e.\ }
\newcommand{\cf}{cf.\ }
\newcommand{\resp}{respectively\ }
\newcommand{\Subsection}{\subsection}
\providecommand{\nobreakdash}{\nobreak\hbox{-}}
\setlength{\emergencystretch}{2em}
\multlinegap0pt
\usepackage{bm}
\usepackage{bbm}
\usepackage{dsfont}
\usepackage{xspace}
\usepackage{cancel}
\usepackage{mathtools}
\usepackage{amsthm}
\makeatletter
\@ifundefined{theorem}{\newtheorem{theorem}{Theorem}}{}
\@ifundefined{lemma}{\newtheorem{lemma}[theorem]{Lemma}}{}
\@ifundefined{proposition}{\newtheorem{proposition}[theorem]{Proposition}}{}
\makeatother
\title[Notes on Chevalley Groups and Root Category I]{Notes on Chevalley Groups and Root Category I}
\author{Buyan Li, Jie Xiao}
\date{Source: arXiv:2505.17805v1 (CC BY 4.0)}
\begin{document}
\maketitle

\noindent\emph{Source and licence.} Notes on Chevalley Groups and Root Category I. Version arXiv:2505.17805v1 (CC BY 4.0), distributed under the Creative Commons Attribution 4.0 International licence, as recorded in the arXiv licence metadata for this exact version and shown on the abstract page https://arxiv.org/abs/2505.17805v1. Full rights notice: licenses/arxiv-2505.17805-CC-BY-4.0.txt.\par\smallskip

\noindent\textit{Editorial scope.} Counted unit: the Proposition whose source label is \texttt{None} in the article (source0.tex lines 2377--2380, source label \texttt{None}), delivered together with the article's own complete original proof of it (source0.tex lines 2382--2482, one \texttt{proof} environment of 101 lines). No mathematics is written, repaired or re-derived: statement and proof are the article's own text line for line. Required context retained uncounted: ABG (lines 804--806). Every definition, display and label of those lines is retained unchanged. Omitted from this package and not redistributed: 1--803, 807--2376, 2381--2381, 2483--3541. Nothing inside a retained excerpt is omitted or altered: every retained body is delivered exactly as the article's lines are written, so no omission receipt of any kind accompanies this package. Citations: the retained excerpts carry no citation command at all, so no bibliography entry of the article is transcribed and no reference is redistributed. Reference adaptation: none was needed and none was made; every reference inside the delivered unit resolves inside it, so no unresolved reference or citation is produced. Printed slips kept literal: no printed word, symbol, hypothesis or number of the counted statement or of its proof was altered. Layout and class: the article's own class is \texttt{amsart} with packages including amsmath, amssymb, latexsym, cancel, rotating, graphicx, mathrsfs, color, fancyhdr, amsthm, xy, verbatim, tikz-cd; the wrapper is the lane's pinned \texttt{amsart} editorial wrapper at 10pt on A4, which restates the theorem-like environments the retained text needs and adds the article's own macro definitions that the retained text needs (none), with extra wrapper packages (bm, bbm, dsfont, xspace, cancel, mathtools). The delivered numbering is the unit's own. Assets: no retained span contains a figure environment, a graphics-inclusion command, a PDF-page inclusion, a table or a TikZ picture; no image file is copied into this package, the package declares no asset dependency and no image file is redistributed with it. Licence: version arXiv:2505.17805v1 of this article is distributed under the Creative Commons Attribution 4.0 International licence, as recorded in the arXiv licence metadata for this exact version and shown on the abstract page https://arxiv.org/abs/2505.17805v1. A full rights notice is delivered at licenses/arxiv-2505.17805-CC-BY-4.0.txt. Characters: every retained excerpt is pure ASCII, so the delivered engine, XeLaTeX with its default Latin Modern text font, reports no missing character.
\begin{lemma}\label{ABG}
    For any $X,Y\in \operatorname{ind}\mathcal{R}$ such that $X\ncong Y$ and $X\ncong TY$, the subcategory of $\mathcal{R}$ generated by $X,Y$ is of type $A_1\times A_1,A_2,B_2$ or $G_2$.
\end{lemma}

\begin{proposition}
    Let $G'$ denote the commutator subgroup of $G$. 
    Then $G=G'$ except for cases $A_1(2),A_1(3),B_2(2),G_2(2)$ ( the notation $A_1(2)$ means $\mathcal{R}$ is of type $A_1$ and the field $\mathbb{K}=\mathbb{F}_2$. Others are similar).
\end{proposition}

\begin{proof}
    Firstly, consider $\mathbb{K}\neq \mathbb{F}_2,\mathbb{F}_3$. 
    Since $h_{[M]}(t) E_{[M]}(u) h_{[M]}(t)^{-1} = E_{[M]}(t^2u)$, we have
    \begin{align*}
        (h_{[M]}(t),E_{[M]}(u)) = E_{[M]}(t^2u-u) \in G'.
    \end{align*}
    Since $\mathbb{K}\neq \mathbb{F}_2,\mathbb{F}_3$, there exists $t\in \mathbb{K}^{\times}$ such that $t^2 - 1 \neq 0$.
    Take $u=\frac{w}{t^2-1}$, and thus $E_{[M]}(w)\in G'$ for any $w\in \mathbb{K}$.
   Since $G'$ contains all the generators of $G$, we have $G=G'$.

   Secondly, consider $\mathbb{K}=\mathbb{F}_3$ and $\mathcal{R}$ is not of type $A_1$.
   For any $M\in \operatorname{ind}\mathcal{R}$, there exists $X\in \operatorname{ind}\mathcal{R}$ such that $M,X$ have extension.
   For example, we may take $X=\tau^{-1}M $, where $\tau$ is the AR-translation functor.
   By Lemma \ref{ABG}, the subcategory $\mathcal{C}$ of $\mathcal{R}$ generated by $M,X$ is of type $A_2,B_2$ or $G_2$.
   We consider these cases respectively.

   \textbf{Case} type $A_2$
   \begin{align*}
    (E_{[TX]}(t),E_{[L_{M,X,1,1}]}(s)) = E_{[M]}(ts\gamma_{TX,L_{M,X,1,1}}^{M} ) \in G'.
   \end{align*} 
   Since $\gamma_{TX,L_{M,X,1,1}}^{M} = \pm 1$, we have $E_{[M]}\subseteq  G'$.

   \textbf{Case} type $B_2$

   In the subcategory $\mathcal{C}$, for all objects $Y$, $d(Y)$ only has two possible values.
   If $d(M)$ is the larger one,
   \begin{align*}
    (E_{[TX]}(t),E_{[L_{M,X,1,1}]}(s)) = E_{[M]}(ts\gamma_{TX,L_{M,X,1,1}}^{M} ) \in G'.
   \end{align*}
   Since $\gamma_{TX,L_{M,X,1,1}}^{M} = \pm 2$, we have $E_{[M]}\subseteq  G'$.

   If $d(M)$ is the smaller one, 
   \begin{align*}
    (E_{[M]}(t),E_{[L_{M,X,1,1}]}(s)) = E_{[L_{M,X,2,1}]}(ts\gamma_{M,L_{M,X,1,1}}^{L_{M,X,2,1}} ) \in G'.
   \end{align*}
   Since $\gamma_{M,L_{M,X,1,1}}^{L_{M,X,2,1}} = \pm 2$, we have $E_{[L_{M,X,2,1}]}\subseteq  G'$.
   Since
   \begin{align*}
    &(E_{[L_{M,X,1,1}]}(a),E_{[TX]}(b)) \\
    =& E_{[M]}(ab\gamma_{L_{M,X,1,1},TX}^{M}) E_{[L_{M,X,2,1}]}(\frac{a^2b}{2!}\gamma_{L_{M,X,1,1},TX}^{M}\gamma_{L_{M,X,1,1},M}^{L_{M,X,2,1}} ) \in G'
   \end{align*}
   and $\gamma_{L_{M,X,1,1},TX}^{M} = \pm 1$, we have $E_{[M]}\subseteq  G'$.

   \textbf{Case} type $G_2$
   \begin{align*}
    (h_{[X]}(t),E_{[M]}(s)) = E_{[M]}((t^{A_{XM}}-1)s) \in G'.
   \end{align*}
   There exists $t$ such that $t^{A_{XM}}-1 \neq 0$, hence $E_{[M]}\subseteq  G'$.

   In all cases, $G'$ contains all the generators of $G$. Hence $G=G'$.

   Finally, consider $\mathbb{K}=\mathbb{F}_2$ and $\mathcal{R}$ is not of type $A_1,B_2,G_2$.
   Since
   \begin{align*}
    (n_{[X]}(t),E_{[Y]}(s)) = E_{[\omega_X(Y)]}(t^{-A_{XY}}s) E_{[Y]}(s)^{-1} \in G',
   \end{align*}
   i.e.
   \begin{align*}
    (n_{[X]}(1),E_{[Y]}(1)) = E_{[\omega_X(Y)]}(1) E_{[Y]}(1)^{-1} \in G',
   \end{align*}
   it suffice to show for a representative of each $W$-orbit of $\operatorname{ind}\mathcal{R}$.
   Notice that $(-|-)$ is $W$-invariant, so does $d(-)$. 
   On the other hand, since for any $M\in \operatorname{ind}\mathcal{R}$, there exists $Y\in \{ S_{1},\cdots,S_{m} \}$ and $i_1,\cdots,i_l$ such that $M\cong \omega_{S_{i_1}}\cdots\omega_{S_{i_l}}(Y)$,
   we only need to show $E_{[S_i]}\subseteq  G'$ for $i=1,\cdots,m$.
   It's easy to see that $d(S_1),\cdots,d(S_m)$ have at most two values.
   Hence $d(\operatorname{ind}\mathcal{R})$ have at most two values.

   We claim that if $d(X) = d(Y), X,Y\in \operatorname{ind}\mathcal{R}$, then $X,Y$ are in the same $W$-orbit.
   We only need to prove for the case when $X,Y$ are simple objects with $(H_X|H_Y)\neq 0$.
   Suppose $S_i,S_j$ are two simple objects satisfying $d(S_i)=d(S_j)$ and $(H_{S_i}|H_{S_j}) \neq 0$.
   Then 
   \begin{align*}
    A_{S_i,S_j} = \frac{(H_{S_i}|H_{S_j})}{d(S_i)} =  \frac{(H_{S_i}|H_{S_j})}{d(S_j)} = A_{S_j,S_i}.
   \end{align*}
   Since $A_{S_i,S_j},A_{S_j,S_i}$ are elements in the Cartan matrix, we must have $A_{S_i,S_j} = A_{S_j,S_i}=-1$.
   (Otherwise they can't be equal.)
   Then $\omega_{S_j}\omega_{S_i}(S_j)=S_i$, and $S_i,S_j$ are in the same $W$-orbit.
   The claim is proved.
   Thus we only need to show one object $M$ for each value of $d(\operatorname{ind}\mathcal{R})$, $E_{[M]}$ is in $G'$.

   If all $d(M)$ are equal, and since $\mathcal{R}$ is not of type $A_1$, there exist $X,Y\in \operatorname{ind}\mathcal{R}$ have extension.
   \begin{align*}
    (E_{[X]}(1),E_{[Y]}(1)) = E_{[L_{X,Y,1,1}]}(1) \in G'.
   \end{align*}
   Hence $G=G'$.

   If there're $d(X),d(Y)$ of two different values, we may assume that $d(X)<d(Y)$ and $X,Y$ have extension.
   Since $\mathcal{R}$ is not of type $G_2$, the subcategory of $\mathcal{R}$ generated by $X,Y$ is of type $B_2$. Hence
   \begin{align}
    (E_{[X]}(1),E_{[Y]}(1)) = E_{[L_{X,Y,1,1}]}(1) E_{[L_{X,Y,2,1}]}(1) \in G',
   \end{align}
   where $d(L_{X,Y,1,1})=d(X)$ and $d(L_{X,Y,2,1})=d(Y)$.
   Since $\mathcal{R}$ is not of type $B_2$, in the quiver corresponding to the complete section, there exist two vertices joined by one edge.
   Suppose these two vertex correspond to $S_i,S_j$. Then
   \begin{align*}
    (E_{[S_i]}(1),E_{[S_j]}(1)) = E_{[L_{S_i,S_j,1,1}]}(1) \in G'.
   \end{align*}
   Hence $E_{[M]}\subseteq  G'$ for all $M\in \operatorname{ind}\mathcal{R}$ such that $d(M) = d(L_{S_i,S_j,1,1})$.
   By formula (5), it also holds for all $M\in \operatorname{ind}\mathcal{R}$ such that $d(M)$ have the different value.
   Thus $G'$ contains all the generators of $G$, and $G=G'$.
\end{proof}

\end{document}
